\documentclass[10pt,a4paper]{amsart}
\usepackage{amsmath,amssymb,amsthm,graphicx}
\usepackage{verbatim}
\newenvironment{reference}{\comment}{\endcomment}
%\newenvironment{reference}{}{}
\newenvironment{slogan}{\comment}{\endcomment}
\newenvironment{history}{\comment}{\endcomment}

% For commutative diagrams we use Xy-pic
\usepackage[all]{xy}

% We use 2cell for 2-commutative diagrams.
\xyoption{2cell}
% Theorem environments.
%
\theoremstyle{plain}
\newtheorem{theorem}[subsection]{Theorem}
\newtheorem{proposition}[subsection]{Proposition}
\newtheorem{lemma}[subsection]{Lemma}

\theoremstyle{definition}
\newtheorem{definition}[subsection]{Definition}
\newtheorem{example}[subsection]{Example}
\newtheorem{exercise}[subsection]{Exercise}
\newtheorem{situation}[subsection]{Situation}

\theoremstyle{remark}
\newtheorem{remark}[subsection]{Remark}
\newtheorem{remarks}[subsection]{Remarks}

\numberwithin{equation}{subsection}

% Macros
%
\def\lim{\mathop{\mathrm{lim}}\nolimits}
\def\colim{\mathop{\mathrm{colim}}\nolimits}
\def\Spec{\mathop{\mathrm{Spec}}}
\def\Hom{\mathop{\mathrm{Hom}}\nolimits}
\def\Ext{\mathop{\mathrm{Ext}}\nolimits}
\def\SheafHom{\mathop{\mathcal{H}\!\mathit{om}}\nolimits}
\def\SheafExt{\mathop{\mathcal{E}\!\mathit{xt}}\nolimits}
\def\Sch{\mathit{Sch}}
\def\Mor{\mathop{\mathrm{Mor}}\nolimits}
\def\Ob{\mathop{\mathrm{Ob}}\nolimits}
\def\Sh{\mathop{\mathit{Sh}}\nolimits}
\def\NL{\mathop{N\!L}\nolimits}
\def\CH{\mathop{\mathrm{CH}}\nolimits}
\def\proetale{{pro\text{-}\acute{e}tale}}
\def\etale{{\acute{e}tale}}
\def\QCoh{\mathit{QCoh}}
\def\Ker{\mathop{\mathrm{Ker}}}
\def\Im{\mathop{\mathrm{Im}}}
\def\Coker{\mathop{\mathrm{Coker}}}
\def\Coim{\mathop{\mathrm{Coim}}}

% Boxtimes
%
\DeclareMathSymbol{\boxtimes}{\mathbin}{AMSa}{"02}

%
% Macros for moduli stacks/spaces
%
\def\QCohstack{\mathcal{QC}\!\mathit{oh}}
\def\Cohstack{\mathcal{C}\!\mathit{oh}}
\def\Spacesstack{\mathcal{S}\!\mathit{paces}}
\def\Quotfunctor{\mathrm{Quot}}
\def\Hilbfunctor{\mathrm{Hilb}}
\def\Curvesstack{\mathcal{C}\!\mathit{urves}}
\def\Polarizedstack{\mathcal{P}\!\mathit{olarized}}
\def\Complexesstack{\mathcal{C}\!\mathit{omplexes}}
% \Pic is the operator that assigns to X its picard group, usage \Pic(X)
% \Picardstack_{X/B} denotes the Picard stack of X over B
% \Picardfunctor_{X/B} denotes the Picard functor of X over B
\def\Pic{\mathop{\mathrm{Pic}}\nolimits}
\def\Picardstack{\mathcal{P}\!\mathit{ic}}
\def\Picardfunctor{\mathrm{Pic}}
\def\Deformationcategory{\mathcal{D}\!\mathit{ef}}
\usepackage{cleveref}
\crefname{originalreference}{original reference}{original references}
\Crefname{originalreference}{Original reference}{Original references}
\makeatletter
\newlabel{lemma-characterize-surjective}{{03H5}{source}{}{sourcepointer-1}{}}
\newlabel{lemma-characterize-surjective@cref}{{[originalreference][0][]03H5}{[1][0][]source}{}{}{}}
\makeatother
\crefname{originalreference}{original reference}{original references}
\Crefname{originalreference}{Original reference}{Original references}
\makeatletter
\newlabel{lemma-points-presentation}{{03BW}{source}{}{sourcepointer-2}{}}
\newlabel{lemma-points-presentation@cref}{{[originalreference][0][]03BW}{[1][0][]source}{}{}{}}
\makeatother
\crefname{originalreference}{original reference}{original references}
\Crefname{originalreference}{Original reference}{Original references}
\makeatletter
\newlabel{lemma-scheme-points}{{03BV}{source}{}{sourcepointer-3}{}}
\newlabel{lemma-scheme-points@cref}{{[originalreference][0][]03BV}{[1][0][]source}{}{}{}}
\makeatother
\crefname{originallemma}{original lemma}{original lemmas}
\Crefname{originallemma}{Original lemma}{Original lemmas}
\makeatletter
\newlabel{morphisms-lemma-etale-open}{{03WT}{source}{}{sourcepointer-4}{}}
\newlabel{morphisms-lemma-etale-open@cref}{{[originallemma][0][]03WT}{[1][0][]source}{}{}{}}
\makeatother
\crefname{originallemma}{original lemma}{original lemmas}
\Crefname{originallemma}{Original lemma}{Original lemmas}
\makeatletter
\newlabel{spaces-lemma-lift-morphism-presentations}{{02X1}{source}{}{sourcepointer-5}{}}
\newlabel{spaces-lemma-lift-morphism-presentations@cref}{{[originallemma][0][]02X1}{[1][0][]source}{}{}{}}
\makeatother
\crefname{originallemma}{original lemma}{original lemmas}
\Crefname{originallemma}{Original lemma}{Original lemmas}
\makeatletter
\newlabel{topology-lemma-quotient}{{08ZK}{source}{}{sourcepointer-6}{}}
\newlabel{topology-lemma-quotient@cref}{{[originallemma][0][]08ZK}{[1][0][]source}{}{}{}}
\makeatother
\title{Existence and uniqueness of the topology on the points of an algebraic space}
\author{The Stacks Project Authors}
\date{stacks-project:spaces-properties.tex:local-snapshot-20261012; GNU FDL 1.2 or later}
\begin{document}
\maketitle
\noindent\textit{Editorial note.} All delivered mathematical text is the author's own, copied from the pinned Stacks Project chapter file spaces-properties.tex; no statement or proof step was rewritten, reordered or removed.\par\medskip
\noindent\textit{Editorial note.} The counted proof is the author's complete argument: the definition of the topology through a presentation, the proof of independence of the presentation through a fibre product square with openness and surjectivity bookkeeping, and the derivation of continuity from a compatible square of presentations.\par\medskip
\noindent\textit{Editorial note.} The label lemma-topology-points names the counted statement, and the proof environment immediately following it in the source is delivered verbatim as its proof.\par\medskip
\noindent\textit{Editorial note.} The counted statement is carried by the author's own lemma environment in the chapter source.\par\medskip
\noindent\textit{Editorial note.} Pre-open gate: python operations/agent-authoring/tools/precard\_gate.py --source spaces-properties.tex --label lemma-topology-points --start 595 --end 674 --local-id 28047 returned RESULT=OK (L28047 free, occupancy FREE, public interval clear). Candidate taken from sources/stacks-spaces-properties/candidates-halfA.json.\par\medskip
\noindent\textit{Editorial note.} The statement and proof are copied from the pinned chapter file; the cited inputs are kept as external pointers with their canonical Stacks tags.\par\medskip
\noindent\textit{Allowed standard prerequisites (editorial).}\begin{itemize}
\item Description of the points of an algebraic space through a presentation. Same chapter, the lemmas identifying the points of an algebraic space with the orbits of the points of a scheme presenting it and characterizing surjectivity on points. Supplies the surjectivity of the map on points from a presentation and from any etale chart.
\item Openness of etale morphisms of schemes on topological spaces. Morphisms, the lemma that an etale morphism of schemes induces an open map on topological spaces. Supplies the openness of the maps appearing in the commutative diagrams of the proof.
\item Lifting a morphism of algebraic spaces to a diagram of presentations. Spaces, the lemma producing a commutative square of presentations for a given morphism of algebraic spaces. Supplies the diagram of schemes over a morphism used to prove continuity.
\end{itemize}\medskip
\ifdefined\hypertarget\hypertarget{sourcepointer-1}{}\fi
\noindent\textit{Source reference.} Original reference 03H5: sources/stacks-spaces-properties/pinned-original-source.tex, source lines 529--529. Prerequisite pointer invoked at proof line(s) 641; it names a standard fact of the same chapter that the counted argument uses, and every step of the counted proof built on it is retained verbatim in the delivered body.\par\medskip
\ifdefined\hypertarget\hypertarget{sourcepointer-2}{}\fi
\noindent\textit{Source reference.} Original reference 03BW: sources/stacks-spaces-properties/pinned-original-source.tex, source lines 569--569. Prerequisite pointer invoked at proof line(s) 643; it names a standard fact of the same chapter that the counted argument uses, and every step of the counted proof built on it is retained verbatim in the delivered body.\par\medskip
\ifdefined\hypertarget\hypertarget{sourcepointer-3}{}\fi
\noindent\textit{Source reference.} Original reference 03BV: sources/stacks-spaces-properties/pinned-original-source.tex, source lines 487--487. Prerequisite pointer invoked at proof line(s) 601; it names a standard fact of the same chapter that the counted argument uses, and every step of the counted proof built on it is retained verbatim in the delivered body.\par\medskip
\ifdefined\hypertarget\hypertarget{sourcepointer-4}{}\fi
\noindent\textit{Source reference.} Original lemma 03WT: sources/stacks-chow/tags.txt, source lines 5124--5124. Prerequisite pointer invoked at proof line(s) 646; it names a standard fact of a sibling chapter that the counted argument uses, and every step of the counted proof built on it is retained verbatim in the delivered body.\par\medskip
\ifdefined\hypertarget\hypertarget{sourcepointer-5}{}\fi
\noindent\textit{Source reference.} Original lemma 02X1: sources/stacks-chow/tags.txt, source lines 3824--3824. Prerequisite pointer invoked at proof line(s) 655; it names a standard fact of a sibling chapter that the counted argument uses, and every step of the counted proof built on it is retained verbatim in the delivered body.\par\medskip
\ifdefined\hypertarget\hypertarget{sourcepointer-6}{}\fi
\noindent\textit{Source reference.} Original lemma 08ZK: sources/stacks-chow/tags.txt, source lines 11378--11378. Prerequisite pointer invoked at proof line(s) 615; it names a standard fact of a sibling chapter that the counted argument uses, and every step of the counted proof built on it is retained verbatim in the delivered body.\par\medskip
\begin{lemma}
\label{lemma-topology-points}
Let $S$ be a scheme. There exists a unique topology on the sets of points
of algebraic spaces over $S$ with the following properties:
\begin{enumerate}
\item if $X$ is a scheme over $S$, then the topology on $|X|$ is the usual one
(via the identification of Lemma \ref{lemma-scheme-points}),
\item for every morphism of algebraic spaces $X \to Y$ over $S$
the map $|X| \to |Y|$ is continuous, and
\item for every \'etale morphism $U \to X$ with $U$ a scheme
the map of topological spaces $|U| \to |X|$ is continuous and open.
\end{enumerate}
\end{lemma}

\begin{proof}
Let $X$ be an algebraic space over $S$. Let $p : U \to X$ be a
surjective \'etale morphism where $U$ is a scheme over $S$.
We define $W \subset |X|$ is open if and only if $|p|^{-1}(W)$
is an open subset of $|U|$. This is a topology on $|X|$
(it is the quotient topology on $|X|$, see
Topology, Lemma \ref{topology-lemma-quotient}).

\medskip\noindent
Let us prove that the topology is independent of the choice of
the presentation. To do this it suffices to show that if $U'$ is a scheme,
and $U' \to X$ is an \'etale morphism, then the map $|U'| \to |X|$
(with topology on $|X|$ defined using $U \to X$ as above)
is open and continuous; which in addition will prove that (3) holds.
Set $U'' = U \times_X U'$, so that we have the commutative diagram
$$
\xymatrix{
U'' \ar[r] \ar[d] & U' \ar[d] \\
U \ar[r] & X
}
$$
As $U \to X$ and $U' \to X$ are \'etale we see that
both $U'' \to U$ and $U'' \to U'$ are \'etale morphisms of schemes.
Moreover, $U'' \to U'$ is surjective. Hence
we get a commutative diagram of maps of sets
$$
\xymatrix{
|U''| \ar[r] \ar[d] & |U'| \ar[d] \\
|U| \ar[r] & |X|
}
$$
The lower horizontal arrow is surjective (see
Lemma \ref{lemma-characterize-surjective}
or
Lemma \ref{lemma-points-presentation})
and continuous by definition of the topology on $|X|$.
The top horizontal arrow is surjective, continuous, and open by
Morphisms, Lemma \ref{morphisms-lemma-etale-open}.
The left vertical arrow is continuous and open (by
Morphisms, Lemma \ref{morphisms-lemma-etale-open}
again.) Hence it follows formally that the right vertical
arrow is continuous and open.

\medskip\noindent
To finish the proof we prove (2).
Let $a : X \to Y$ be a morphism of algebraic spaces. According to
Spaces, Lemma \ref{spaces-lemma-lift-morphism-presentations}
we can find a diagram
$$
\xymatrix{
U \ar[d]_p \ar[r]_\alpha & V \ar[d]^q \\
X \ar[r]^a & Y
}
$$
where $U$ and $V$ are schemes, and $p$ and $q$ are surjective and \'etale.
This gives rise to the diagram
$$
\xymatrix{
|U| \ar[d]_p \ar[r]_\alpha & |V| \ar[d]^q \\
|X| \ar[r]^a & |Y|
}
$$
where all but the lower horizontal arrows are known to be continuous and
the two vertical arrows are surjective and open. It follows that the
lower horizontal arrow is continuous as desired.
\end{proof}

\end{document}
